% BHU PYQ -- Transcribed & Corrected
% Paper: PHYSE31: Computational Physics: Programming and Data Analysis
% Exam: Under Graduate Semester III (NEP) Examination, 2025-26
% Ref Code: CECONF/N/2025-26/007
% Category: Skill Enhancement Course (SEC)
% Department: Physics

\documentclass[11pt,a4paper]{article}
\usepackage[a4paper,top=2.5cm,bottom=2.5cm,left=2.8cm,right=2.8cm,headheight=14pt]{geometry}
\usepackage{amsmath,amssymb}
\usepackage[T1]{fontenc}
\usepackage[utf8]{inputenc}
\usepackage{lmodern,microtype,fancyhdr,enumitem,hyperref,lastpage,parskip,listings}

\hypersetup{
    colorlinks=true,
    urlcolor=black,
    linkcolor=black,
    pdftitle={PHYSE31: Computational Physics (End Term) -- BHU Sem III 2025-26 (NEP SEC)}
}

\pagestyle{fancy}
\fancyhf{}
\renewcommand{\headrulewidth}{0.4pt}
\renewcommand{\footrulewidth}{0.4pt}
\fancyhead[L]{\small \textit{Banaras Hindu University}}
\fancyhead[R]{\small \textit{PHYSE31: Computational Physics (2025-26)}}
\fancyfoot[C]{\small Page \thepage\ of \pageref{LastPage}}

\newcommand{\pts}[1]{\hfill \textit{[#1]}}

\lstset{
  language=Python,
  basicstyle=\ttfamily\small,
  keywordstyle=\bfseries,
  showstringspaces=false,
  frame=single,
  frameround=tttt
}

\begin{document}

\begin{center}
  {\large\bfseries BANARAS HINDU UNIVERSITY}\\[2pt]
  {\normalsize Under Graduate Semester III Examination, 2025-26 (NEP)}\\[6pt]
  \rule{0.8\linewidth}{0.5pt}\\[6pt]
  {\large\bfseries PHYSICS (SEC)}\\[4pt]
  {\large\bfseries Paper Code: PHYSE31 : Computational Physics : Programming and Data Analysis}\\[2pt]
  {\small Ref. No.: CECONF/N/2025-26/007}\\[4pt]
  {\normalsize Time: 2$\frac{1}{2}$ Hours\hfill Maximum Marks: 70}
\end{center}

\rule{\linewidth}{0.4pt}

\smallskip

\noindent \textit{Instructions: Answer five questions including Question 1 which is compulsory. Use Python programming language to address all the questions.}

\bigskip

\noindent \textbf{Question 1.} Answer \textit{any four} of the following questions: \pts{4 $\times$ 5 = 20}
\begin{enumerate}[label=(\alph*)]
  \item Program to print the multiplication table of 4 using `for` loop.
  \item Program to count digits in your Exam roll number.
  \item Program to reverse the statement ``We Are Happy''.
  \item Write the value that gets printed, or write ``error'' if there is an error during the execution:
\begin{lstlisting}
numbers = [4, 2, -1]
x = 1
for num in numbers:
    x = x * numbers
print(x)
\end{lstlisting}
  \item Write the value that gets printed, or write ``error'' if there is an error during the execution:
\begin{lstlisting}
x = [1, 2, 3]
x[len(x) - 1] = x[len(x) - x[0]] - x[1]
print(x[2])
\end{lstlisting}
\end{enumerate}

\bigskip

\noindent \textbf{Question 2.}
\begin{enumerate}[label=(\alph*)]
  \item Write a program to print: \pts{6}
  \begin{enumerate}[label=(\roman*)]
    \item your semester and course name.
    \item your enrolment number.
    \item name of your examination.
  \end{enumerate}
  \item Find the data type of: \pts{6$\frac{1}{2}$}
  \begin{enumerate}[label=(\roman*)]
    \item your course name.
    \item your enrolment number.
    \item name of your examination.
  \end{enumerate}
\end{enumerate}

\bigskip

\noindent \textbf{Question 3.}
\begin{enumerate}[label=(\alph*)]
  \item Write a python code to find the biggest of three numbers 'x', 'y' and 'z'. Thereafter, assign the values of your choice to the variables and obtain the results. \pts{8}
  \item Explain `continue` and `break` statements with an example. \pts{4$\frac{1}{2}$}
\end{enumerate}

\bigskip

\noindent \textbf{Question 4.}
\begin{enumerate}[label=(\alph*)]
  \item Given \texttt{fruits = ["apple", "banana", "cherry", "date"]}: \pts{6$\frac{1}{2}$}
  \begin{enumerate}[label=(\roman*)]
    \item Change the first element of the above list to orange.
    \item Change the above list to tuple.
    \item What is the difference between a list and a tuple?
    \item Print each fruit in the list 'fruits'.
  \end{enumerate}
  \item Discuss the following list functions with the help of examples: \pts{6}
  \begin{enumerate}[label=(\roman*)]
    \item \texttt{len()}
    \item \texttt{sum()}
    \item \texttt{any()}
  \end{enumerate}
\end{enumerate}

\bigskip

\noindent \textbf{Question 5.} Make a $4 \times 4$ array (of integral numbers) using \texttt{numpy}. \pts{12$\frac{1}{2}$}
\begin{enumerate}[label=(\alph*)]
  \item Slice a $2 \times 2$ submatrix out of it.
  \item Find the mean of all the rows.
  \item Find the standard deviation of all the columns.
  \item Write these outputs to a file named 'int\_mat.txt'.
\end{enumerate}

\bigskip

\noindent \textbf{Question 6.} Given: \pts{12$\frac{1}{2}$}
\begin{lstlisting}
Name   : ['Sam', 'Andrea', 'Alex', 'Robin', 'Kia']
Age    : [14, 25, 55, 8, 21]
Gender : ['M', 'F', 'M', 'M', 'F']
Weight : [45, 88, 56, 15, 71]
\end{lstlisting}
Use the above data to create a data frame and calculate the following:
\begin{enumerate}[label=(\alph*)]
  \item Sort the dataframe by age in descending order.
  \item Divide the dataframe based on gender.
  \item What are three key differences between numpy and pandas?
\end{enumerate}

\bigskip

\noindent \textbf{Question 7.} Given $x = (1, 2, 3, 4, 5, 6, 7)$, write a python program to draw the curves between $x$ vs $x^2$ and $x$ vs $x^3$ in a single figure. Add a title, axes label and legends on the plot. How do you change the legend fontsize in matplotlib? \pts{12$\frac{1}{2}$}

\bigskip
\begin{center}
  \textbf{***}
\end{center}

\end{document}
